\documentclass[xcolor=x11names,compress,aspectratio=169]{beamer}
%\documentclass[xcolor=x11names,compress,aspectratio=43]{beamer}

%% General document %%%%%%%%%%%%%%%%%%%%%%%%%%%%%%%%%%
\usepackage{graphicx}
\usepackage{tikz}
\usepackage{amsmath}
\usepackage{amssymb}
\usepackage[utf8]{inputenc}
\usepackage{textpos}
\usepackage{booktabs}
\usepackage{listings}
\usepackage{auto-pst-pdf}

\usetikzlibrary{decorations.fractals, calc}
%%%%%%%%%%%%%%%%%%%%%%%%%%%%%%%%%%%%%%%%%%%%%%%%%%%%%%

%% Beamer Layout %%%%%%%%%%%%%%%%%%%%%%%%%%%%%%%%%%
\useoutertheme[subsection=false,shadow]{miniframes}
\useinnertheme{default}
%\usefonttheme{serif}
\usepackage{palatino}
\setbeamerfont{title like}{shape=\scshape}
\setbeamerfont{frametitle}{shape=\scshape}
\beamertemplatenavigationsymbolsempty


% Uncomment this line, if you want frame numbers on your slides
% (I disabled them since I had the progress bar at the header)
%\setbeamertemplate{footline}[frame number]

\definecolor{fmiBlue}{RGB}{11,128,145} % FSU Faculty blue
\definecolor{evbc}{RGB}{50,167,132} % EVBC Green
\definecolor{bgGray}{RGB}{230,230,230} % some gray if needed


\setbeamercolor*{upper separation line head}{bg=Snow4} %Snow4 comes from the x11names option of xcolor
%\setbeamercolor*{lower separation line head}{bg=evbc}  % obsolete - there is now a cool progress bar 
\setbeamercolor*{normal text}{fg=black,bg=white} 
\setbeamercolor*{alerted text}{fg=red} 
\setbeamercolor*{example text}{fg=black} 
\setbeamercolor*{structure}{fg=black} 
 
\setbeamercolor*{palette tertiary}{fg=black,bg=black!10} 
\setbeamercolor*{palette quaternary}{fg=black,bg=black!10} 

%%%%%%%%%%%%%%%%%%%%%%%%%%%%%%%%%%%%%%%%%%%%
% IMPORTANT
% Kevin:
% Currently the FMI blue (the light uni blue) is here.
% If you want for example the evbc green, you'd have to
% change the colors here. I will put this in a nice function
% at some point.
\setbeamercolor{title}{fg=evbc}
\setbeamercolor{structure}{fg=evbc}
\setbeamercolor{frametitle}{fg=evbc}
\setbeamercolor{block title}{fg=evbc, bg=bgGray}
\setbeamercolor{block title example}{fg=evbc}
\setbeamercolor{enumerate item}{fg=evbc}
\setbeamercolor{itemize item}{fg=evbc}
\setbeamercolor{item projected}{bg=evbc}

\setbeamerfont{block title}{size=\small}
\setbeamerfont{block body}{size=\scriptsize}

% Use those if you want. Just some shortcuts for 
% columns. Personally, I use minipages nowadays.
\renewcommand{\(}{\begin{columns}}
\renewcommand{\)}{\end{columns}}
\newcommand{\<}[1]{\begin{column}{#1}}
\renewcommand{\>}{\end{column}}

% Kevin:
% This is used for the code examples in the document.
% If used, set the frame option to 'fragile' (!!!)
\lstset{
	language=sh, 
	numbers=left, 
	numberstyle=\tiny\color{gray}, 
	backgroundcolor=\color{bgGray}, 
	basicstyle=\scriptsize\ttfamily,
	commentstyle=\tiny\color{Green4},
	keywordstyle=\color{Blue3},
	escapeinside=@@
}

% Kevin
% Appendix page number fix. More like an ugly hack.
% I was never able to create a "nice" appendix command, however, this one here works fine.
\newcommand{\beginbackup}{
   \newcounter{framenumbervorappendix}
   \setcounter{framenumbervorappendix}{\value{framenumber}}
}
\newcommand{\backupend}{
   \addtocounter{framenumbervorappendix}{-\value{framenumber}}
   \addtocounter{framenumber}{\value{framenumbervorappendix}} 
}

%%%%%%%%%%%%%%%%%%%%%%%%%%%%%%%%%%%%%%%%%%%%%%%%%%
% Kevin: Change this, if you want.
\graphicspath{{figures/}}
%

% Kevin:
% I used this slides for subsections a long time ago, however, at some point I figured that it is nicer to have those slides
% at the beginning of each section. Use this extra slides to give a short summary of the previous talk.
\AtBeginSection[]{
  \begin{frame}
  \vfill
  \centering
  \begin{beamercolorbox}[sep=8pt,center,shadow=true,rounded=true]{title}
    \usebeamerfont{title}\insertsection\par%
  \end{beamercolorbox}
  \vfill
  \end{frame}
}


% Kevin:
% small macro, used to build the progress bar
\def\insertframeratio{%
    \pgfmathparse{\insertframenumber/\inserttotalframenumber}%
}
%%%%%%%%%%%%%%%%%%%%%%%%%%
% Kevin:
% This is the progress bar in the headline.
% Change the textcolor if needed, change the width to 2pt if wanted
\addtobeamertemplate{headline}{}{%
  \insertframeratio
  \edef\myvar{\pgfmathresult}
  %\textcolor{evbc}{\rule{\myvar \paperwidth}{1pt}}
  \textcolor{evbc}{\rule{\paperwidth}{1pt}}
}

%%%%%%%%%%%%%%%%%%%%%%%%%%%
% Kevin
% This adds the FSU Logo to the
% bottom left. Changing the logo is easy
% adding a second logo (e.g. the evbc) is tricky
\addtobeamertemplate{footline}{}{%
\begin{textblock*}{100mm}(1em,-3.5em)
\includegraphics[scale=0.4]{evbc_cmyk.pdf}
\end{textblock*} 
\begin{textblock*}{100mm}(73em,-1em) 
\insertframenumber/\inserttotalframenumber
\end{textblock*}
}

%% Kevin:
%% I used this in order to show two logos at the same time.
%% as you can see, it is quite tricky in terms of positioning 
%% and size, but a nice try-and-error session usually
%% works out.
%\addtobeamertemplate{footline}{}{%
%\begin{textblock*}{100mm}(0.5em,-5em)
%\includegraphics[scale=0.5]{figures/evbc_cmyk.pdf}
%\end{textblock*}
%\begin{textblock*}{100mm}(11em,-3em)
%\includegraphics[scale=0.025]{fsu_logo.jpg}
%\end{textblock*}
%}


% Kevin:
% Well, fill this out as needed
%\title{Theoretical and practical metagenomic approaches to viral discovery}
%\subtitle{}
%\author{Kevin Lamkiewicz, Manja Marz}
%\date{27.09.2018\\[1em]RNA Bioinformatics and High-Throughput Analysis\\[1em]Friedrich Schiller University Jena}
%\date{xx.xx.2018\\RNA Bioinformatics and High-Throughput Analysis\\[2em]\includegraphics[scale=0.05]{fsu_logo.jpg}}



\title{Theoretical and practical metagenomic approaches to viral discovery}
\subtitle{Practical Session: Practical Introduction to Machine Learning and SVMs}
\author{Kevin Lamkiewicz, Manja Marz}
\date{24.10.2019\\[1em]European Virus Bioinformatics Center}

\begin{document}

\begin{frame}
  \maketitle
\end{frame}

\begin{frame}[c]\frametitle{whoami}
  \uncover<2->{
  \begin{minipage}{0.5\textwidth}
    \begin{itemize}
      \item PhD student at Manja's Lab since 2017
      \item Interested in (RNA) viruses 
      \item Focus on RNA secondary structures and ncRNAs
      \item Experience with Machine Learning, RNA structures, assemblies, alignments, SNP calling, \dots
    \end{itemize}
  \end{minipage} \hfill \begin{minipage}{0.45\textwidth}
    \begin{figure}
      \centering
      \includegraphics[width=0.5\textwidth]{strahlemann.jpg}
    \end{figure}
  \end{minipage}
  }
\end{frame}

\section[sklearn]{Introducing scikit-learn}

\begin{frame}[c,fragile]\frametitle{Building a support vector machine with Python3}
  \begin{lstlisting}[language=Python]
#!/usr/bin/env python3

# to import scikit-learn we write the following
import sklearn
@\pause@
# but actually we need this:
from sklearn import svm
@\pause@
# now we create our SVM
classifier = svm.SVC(kernel='linear')
@\pause@
# Done.
  \end{lstlisting}
\end{frame}

\begin{frame}[c,fragile]\frametitle{Understanding the two lines of code}

  \begin{lstlisting}[language=Python]
from sklearn import svm
  \end{lstlisting}
  \begin{itemize}
    \item {\small make all functions of the \texttt{svm} class available}
  \end{itemize}
  
\pause
  \begin{lstlisting}[language=Python]
classifier = svm.SVC(kernel='linear')
  \end{lstlisting}
  \begin{itemize}
    \item {\small create a classifier object with a linear kernel}
  \end{itemize}

\end{frame}

\section[Iris data]{Applying machine learning to real biological data}

\begin{frame}[c, fragile]\frametitle{The Iris test data set}

  \begin{lstlisting}[language=Python]
from sklearn import svm, datasets
from sklearn.metrics import confusion_matrix, accuracy_score
# load the iris dataset
iris = datasets.load_iris()
@\pause@
# parse the important fields
data = iris.data
target = iris.target
@\pause@
# create a classifer model via svm and train it with the data
svmLinear = svm.SVC(kernel='linear').fit(data, target)
@\pause@
# predict with our model
prediction = svmLinear.predict(data)
@\pause@
# check performance of our prediction and model
confusion_matrix(target, prediction)
accuracy_score(target, prediction)
  \end{lstlisting}
\end{frame}

\begin{frame}[c, fragile]\frametitle{Easy, right?}
  \begin{block}{Problems...!?}
    \uncover<2->{
      I used the same data for training and prediction (testing). This is very bad practice and 
      should be avoided at all cost.\\ \ \\
      Our classifier learns based on the given training data - obviously, if the test data is identical
      to the training data, the model performs well.\\ \ \\
      We have to split our training set...
    }
  \end{block}
\end{frame}

\begin{frame}[c, fragile]\frametitle{Splitting the data in training and test sets}
  \begin{lstlisting}[language=Python]
from sklearn.model_selection import train_test_split
@\pause@
# here, we split our data set into different parts - for training and testing
# the method train_test_split() returns a tuple with 4 elements,
# which are stored in the appropriate variables, respectively.
data_train, data_test, target_train, target_test = train_test_split(
        data, target, test_size=0.33)
@\pause@
# train the model on the split data
svmLinear = svm.SVC(kernel='linear').fit(data_train, target_train)
# predict the targets of the "new data"
prediction = svmLinear.predict(data_test)
@\pause@
# get the performance of this model
confusion_matrix(target_test, prediction)
accuracy_score(target_test, prediction)
  \end{lstlisting}
\end{frame}

\section[Hyperparameter]{SVM Kernel and other Hyperparamter}
\begin{frame}[c, fragile]\frametitle{Different kernels may yield different results...}
  \uncover<2->{
  \begin{figure}
    \centering
    \includegraphics[width=0.6\textwidth]{sphx_glr_plot_iris_001.png}
  \end{figure}
  }
\end{frame}

\begin{frame}[c, fragile]{Let us classify handwritten digits}
  \begin{lstlisting}[language=Python]
from sklearn import svm, datasets
from sklearn.metrics import confusion_matrix, accuracy_score
from sklearn.model_selection import train_test_split
# get the data
digits = datasets.load_digits()
@\pause@
# split everything
data_train, data_test, target_train, target_test = train_test_split(
        data, target, test_size=0.33)
# train the model with RBF kernel on the training set
digitSVM = svm.SVC(kernel='rbf').fit(data_train, target_train)
@\pause@
prediction = digitSVM.predict(data_test)
confusion_matrix(target_test, prediction)
accuracy_score(target_test, prediction)
  \end{lstlisting}
\end{frame}

\begin{frame}[c, fragile]{Well, that didn't work...}
  \pause
  \begin{lstlisting}[language=Python]
# train the model with linear kernel on the training set
digitSVM = svm.SVC(kernel='linear').fit(data_train, target_train)
@\pause@
prediction = digitSVM.predict(data_test)
confusion_matrix(target_test, prediction)
accuracy_score(target_test, prediction)
  \end{lstlisting}
\end{frame}

\begin{frame}[c,fragile]{Changing the degree of our kernel function}
  A linear kernel is essentially the same as a polynomial kernel with
  degree 1. So, what happens, if we change this parameter?
  \pause
  \begin{lstlisting}[language=Python, showstringspaces=false]
# init a for loop to change the hyperparamter 'degree'
for d in range(1,11):
  # train different SVMs with a different kernel each time
  digitSVM = svm.SVC(kernel='poly', degree=d, gamma='auto').
    fit(data_train, target_train)
  prediction = digitSVM.predict(data_test)
  # some output things
  print(f"Degree of polynomial kernel: {d}")
  print(f"Accuracy of digit classification: 
    {accuracy_score(target_test, prediction)}")
  \end{lstlisting}
\end{frame}

% Backup Slides. Using this macro, you'd get a slide number for each
% backup slide without increasing the maximum slide numbers of the original presentation.
% However, for this, the framenumber has to be inserted - which isn't in the template by default
\beginbackup

\begin{frame}[c]\frametitle{Coffee Break}
  \begin{figure}[htbp]
    \centering
    \includegraphics[width=0.65\textwidth]{coffeebreak.png}
  \end{figure}
\end{frame}

\backupend

\end{document}